\documentclass[10pt,a4paper]{article}
\usepackage{fontspec,amsmath,amssymb,graphicx}
\usepackage[margin=18mm]{geometry}
\setmainfont{Times New Roman}
\setlength{\parindent}{0pt}
\setlength{\parskip}{5pt}
\allowdisplaybreaks
\sloppy
\emergencystretch=2em
\begin{document}
\section*{BT\_5\_5\_3\_11 — Lời giải bổ sung}
Nếu $u$ hằng trên một đĩa mở, tính giải tích và liên thông suy hằng trên toàn $D(1)$, trái giả thiết. Với $0<r_1<r_2$, cực đại trên $\overline D(r_1)$ nhỏ nghiêm ngặt cực đại trên $\overline D(r_2)$: nếu bằng nhau, $u$ đạt cực đại ở điểm nội của đĩa lớn nên hằng. Với $r_1=0$, dùng cực đại mạnh tương tự. Vì vậy $\phi$ tăng nghiêm ngặt, kéo theo không giảm.\end{document}